\section{Actual original complex Taylor polynomials and stable coefficients}

This is the convergent original-coefficient bridge toward the formal
standard-basis setting of \cite[\S7]{bierstone-milman97-invent-math}.
Every polynomial is evaluated at the ACTUAL original coordinate germs,
with the ORIGINAL constant-germ complex action. The already proved
zero-tangential-block identification $\mathbb C^{0+n}=\mathbb C^n$
retains these same coordinates. No formal-series multiplication or division
is asserted in this chapter.
Write $R=\mathcal O_{\mathbb C^n,0}$, $\mathfrak m$ for its actual
maximal ideal, $X^\alpha$ for the actual coordinate monomial, and
$|\alpha|=\sum_i\alpha_i$ for the true original exponent degree.

\begin{lemma}[Actual original polynomial evaluation separates all lower coefficients]
\label{hir:actual-original-polynomial-low-coefficient-separation}
\uses{hir:relative-normal-basis-actual-coefficient-separation,hir:relative-normal-actual-zero-base-complex-algebra}
\lean{Hironaka.HironakaActualOriginalTaylorPolynomial.originalPolynomialEval,Hironaka.HironakaActualOriginalTaylorPolynomial.originalPolynomialEval_monomial,Hironaka.HironakaActualOriginalTaylorPolynomial.normalPolynomialEval_original_coefficient_lift,Hironaka.HironakaActualOriginalTaylorPolynomial.originalPolynomial_coefficient_eq_zero_of_eval_mem}
\leanok
The actual evaluation $\operatorname{ev}:\mathbb C[X_1,\ldots,X_n]\to R$
is an algebra homomorphism with
$\operatorname{ev}(aX^\alpha)=a\,X^\alpha$ in the original whole ring.
If $\operatorname{ev}(P)\in\mathfrak m^q$, EVERY actual coefficient
$[X^\alpha]P$ with $|\alpha|<q$ is zero.
\end{lemma}
\begin{proof}
Use literal original coordinate germs in polynomial evaluation; its
monomial formula is the original finite product formula. Map the same
complex polynomial coefficients to their actual zero-dimensional
constant germs. True zero-block pullback of each such constant is the
SAME original ambient complex constant, so evaluating this coefficient
lift by the existing whole normal polynomial map gives precisely
$\operatorname{ev}(P)$. The actual zero-tangent normal ideal is
$\mathfrak m$. Existing genuine whole normal coefficient separation
therefore annihilates every coefficient below $q$ as a zero-dimensional
whole germ; actual evaluation at its unique point gives the displayed
original complex coefficient zero. This includes zero dimension and
$q=0$, where the lower-coefficient condition is empty.
\end{proof}

\begin{theorem}[Every original numerator has a canonical genuine finite Taylor polynomial]
\label{hir:actual-original-canonical-finite-taylor-polynomial}
\uses{hir:actual-original-polynomial-low-coefficient-separation,hir:relative-normal-actual-base-jet-complex-monomial-dictionary}
\lean{Hironaka.HironakaActualOriginalTaylorPolynomial.originalJetCoordinates,Hironaka.HironakaActualOriginalTaylorPolynomial.originalTaylorPolynomial,Hironaka.HironakaActualOriginalTaylorPolynomial.originalTaylorPolynomial_coefficient_low,Hironaka.HironakaActualOriginalTaylorPolynomial.originalTaylorPolynomial_coefficient_high,Hironaka.HironakaActualOriginalTaylorPolynomial.originalTaylorPolynomial_quotient_class,Hironaka.HironakaActualOriginalTaylorPolynomial.originalTaylorPolynomial_remainder,Hironaka.HironakaActualOriginalTaylorPolynomial.originalTaylorPolynomial_unique,Hironaka.HironakaActualOriginalTaylorPolynomial.originalTaylorPolynomial_add,Hironaka.HironakaActualOriginalTaylorPolynomial.originalTaylorPolynomial_smul}
\leanok
For every original $g\in R$ and $k\ge0$, internally define
$T_k(g)\in\mathbb C[X]$ using the actual distinct-monomial complex
coordinates of the TRUE quotient $R/\mathfrak m^{k+1}$.
Its coefficients vanish when $|\alpha|\ge k+1$, and
\[
 g-\operatorname{ev}(T_k(g))\in\mathfrak m^{k+1}.
\]
It is the UNIQUE polynomial with that coefficient support bound and that
actual whole remainder. The construction is complex linear in the
original numerator.
\end{theorem}
\begin{proof}
Apply the inverse of the already constructed original complex monomial
quotient equivalence to the ACTUAL class of $g$. If its coordinate tuple
is $c_\alpha$, set $T_k(g)=\sum_{|\alpha|<k+1}c_\alpha X^\alpha$.
The genuine finite original exponent family gives the asserted support
bound. Expanding its coefficient tuple as the finite sum of standard
basis vectors and using the actual monomial-class formula for that SAME
quotient equivalence proves that the class of its evaluation is exactly
the original class of $g$. This is precisely the stated whole remainder.
If $P$ has the same support bound and remainder, evaluation of
$P-T_k(g)$ lies in $\mathfrak m^{k+1}$. The first lemma annihilates all
its lower coefficients; both original support bounds annihilate the
others. Polynomial coefficient extensionality gives uniqueness.
All constituent quotient, inverse-basis and coefficient combination maps
are linear for the original complex action, giving linearity.
\end{proof}

\begin{theorem}[True original Taylor coefficients stabilize across all actual truncations]
\label{hir:actual-original-taylor-coefficient-stabilization}
\uses{hir:actual-original-canonical-finite-taylor-polynomial,hir:actual-original-polynomial-low-coefficient-separation}
\lean{Hironaka.HironakaActualOriginalTaylorPolynomial.originalTaylorPolynomial_coefficient_stable,Hironaka.HironakaActualOriginalTaylorPolynomial.originalTaylorCoefficientMap,Hironaka.HironakaActualOriginalTaylorPolynomial.originalTaylorCoefficientMap_eq_coefficient}
\leanok
If $k\le l$ and $|\alpha|<k+1$, then
$[X^\alpha]T_l(g)=[X^\alpha]T_k(g)$.
Thus the original complex coefficient
$c_\alpha(g):=[X^\alpha]T_{|\alpha|}(g)$ equals
$[X^\alpha]T_k(g)$ for EVERY $k$ with $|\alpha|<k+1$.
\end{theorem}
\begin{proof}
The true higher remainder belongs to $\mathfrak m^{l+1}\subseteq
\mathfrak m^{k+1}$, and the true lower remainder belongs there too.
Subtracting these genuine original whole remainders shows that
$\operatorname{ev}(T_l(g)-T_k(g))\in\mathfrak m^{k+1}$.
Actual lower polynomial coefficient separation gives the first equality.
Apply it with lower truncation $|\alpha|$ to obtain the stated equality
for the actual stable coefficient. No coefficient compatibility premise
is supplied by a caller.
\end{proof}

\begin{theorem}[All genuine original coefficients give a faithful complex linear map]
\label{hir:actual-original-all-taylor-coefficients-faithful}
\uses{hir:actual-original-taylor-coefficient-stabilization,hir:actual-whole-ideal-maximal-filtration-separated}
\lean{Hironaka.HironakaActualOriginalTaylorPolynomial.originalTaylorCoefficientMap_eq_zero_iff,Hironaka.HironakaActualOriginalTaylorPolynomial.originalTaylorCoefficientMap_injective}
\leanok
The actual original coefficient map
$g\mapsto(c_\alpha(g))_\alpha$ is complex linear and injective into the
actual coefficient function space $(\mathbb N^n\to\mathbb C)$.
\end{theorem}
\begin{proof}
Linearity follows from the true finite polynomial construction at the
original degree $|\alpha|$. If all actual coefficients of $g$ vanish,
stabilization annihilates every lower coefficient of EVERY $T_k(g)$;
its genuine support bound annihilates the others, so $T_k(g)=0$.
The true remainder then places the original $g$ in EVERY
$\mathfrak m^{k+1}$. Whole original ideal separation at $I=0$ gives
$g=0$. Apply this kernel conclusion to a difference, using actual
complex linearity, to obtain injectivity. This is a faithful original
linear coefficient map; formal-series algebra multiplication remains the
next explicit obligation.
\end{proof}
